\paragraph{Freeness at ramified coefficient primes.}
The missing integral condition is stronger than a characteristic-polynomial
identity. It can be proved on geometric fibres by using the universal
extension of the abelian quotient, followed by a three-step filtration.
The argument uses the full ring of integers, including its nilpotents
after reduction.

\begin{lemma}\label{geo-witt-coefficient-freeness}
Let $k$ be an algebraically closed field of characteristic $p>0$,
$W=W(k)$, and $\mathcal O$ the ring of integers of a number field $E$.
If a finite free $W$-module $D$ has an $\mathcal O$-action, then $D/pD$
is projective over $C=k\otimes_{\mathbf Z}\mathcal O$.
If moreover, for every $a\in\mathcal O$,
\begin{equation}\label{geo-regular-character-freeness}
 \det(X-a\mid D/pD)
 =\det(X-a\mid C)^r,
\end{equation}
for an integer $r\geq0$, then $D/pD$ is free of rank $r$ over $C$.
\end{lemma}
\begin{proof}
The ring $W\otimes_{\mathbf Z}\mathcal O$ is a product of discrete
valuation rings. Indeed, decompose $\mathcal O\otimes\mathbf Z_p$ by
its primes over $p$. The maximal unramified subextension of each local
field splits over $W$, and a uniformizer of the remaining totally
ramified extension has an Eisenstein polynomial. Each resulting factor
is a discrete valuation ring with uniformizer $\varpi$ and
$p=u\varpi^e$, where $u$ is a unit.
The corresponding summand of $D$ has no $\varpi$-torsion, since it has
no $p$-torsion. It is therefore free over this factor. Reduction modulo
$p$ proves projectivity over $C$.

Write $C=\prod_i C_i$, where $C_i\simeq k[t]/(t^{e_i})$ and the
residue character is $\tau_i:\mathcal O\to k$. If the rank on $C_i$
is $r_i$, the characteristic polynomial of $a$ is
$\prod_i(X-\tau_i(a))^{e_i r_i}$.
For a selected prime of $\mathcal O$ above $p$, the Chinese remainder
theorem gives an $a$ which is zero at the other primes and is a
nonzero generator of its residue field at that prime. Its distinct
residue embeddings give distinct nonzero roots. Comparing their
multiplicities in \eqref{geo-regular-character-freeness} gives
$e_i r_i=e_i r$, hence $r_i=r$ for every factor.
\end{proof}

For the next two propositions, retain an algebraically closed field
$k$ of characteristic $p>0$, the full ring of integers $\mathcal O$
of the number field $E$, and $C=k\otimes_{\mathbf Z}\mathcal O$.

\begin{proposition}\label{geo-proper-formal-hodge-comparison}
Let $B/k$ be an abelian variety with an $\mathcal O$-action and a
rigidified line bundle $M$ whose associated map $B\to B^\vee$ is an
isomorphism. Suppose its adjunction identifies the endomorphism
$a$ with the dual of $\bar a$, for an involution
$a\mapsto\bar a$ of $\mathcal O$.
Put $V=\operatorname{Lie}B$.
Choose formal coordinates at zero and a frame $e(x,y)$, rigidified
on both axes, of the associated Poincar\'e biextension.
Write its adjunction in this frame as
\[
             e(ax,y)\longmapsto u_a(x,y)e(x,\bar a y).
\]
If $\rho(a)$ is the tangent action and
$b_a$ is the mixed bilinear coefficient of $u_a$, then the module
$V\oplus V^*$ with action
\begin{equation}\label{geo-formal-hodge-action}
       a(v,\lambda)=
       \bigl(\rho(a)v,\rho(\bar a)^t\lambda-b_a^t v\bigr)
\end{equation}
is projective over $C=k\otimes_{\mathbf Z}\mathcal O$.
\end{proposition}
\begin{proof}
The universal vector extension of $B^\vee$ parametrizes extensions
of $B$ by $\mathbf G_m$ together with a pointed section over the
first infinitesimal neighbourhood of zero. This description is
functorial and commutes with arbitrary base change
\cite[I, \S2.1 and Proposition 2.6.7]{mazur-messing-universal}.
Use $M$ to identify the parameter space $B^\vee$ with $B$.
Over $k[\epsilon]/(\epsilon^2)$, a tangent point above $\epsilon v$
has a unique rigidification of the form
\[
                  e(\epsilon v,y)(1+\epsilon\lambda(y)),
                 \qquad \lambda\in V^*,
\]
on that first neighbourhood. This identifies the tangent space
of the universal extension with $V\oplus V^*$.
Pulling the rigidification back by $\bar a$ and then through the
inverse adjunction gives
\[
 e(a\epsilon v,y)
 \bigl(1+\epsilon\{\rho(\bar a)^t\lambda-b_a^t v\}(y)\bigr).
\]
This proves precisely \eqref{geo-formal-hodge-action}, with its
minus sign. It uses only first infinitesimal neighbourhoods, so
no logarithm or division by $p$ occurs.

The Lie algebra of the universal extension is the reduction of
the finite free $W(k)$-module furnished by its Dieudonn\'e crystal;
the comparison is functorial in the abelian variety and its
endomorphisms
\cite[II, \S9, \S13 and \S15]{mazur-messing-universal}.
Consequently the $\mathcal O$-action lifts to this module.
Lemma~\ref{geo-witt-coefficient-freeness} applies. In particular,
neither the abelian variety nor its endomorphisms need to lift
together to characteristic zero.
\end{proof}

\begin{proposition}\label{geo-toric-hodge-filtration}
Let $G/k$ be a semiabelian variety with an $\mathcal O$-action,
and let $q:G\to B$ be an equivariant quotient with torus kernel.
Put $V=\operatorname{Lie}G$, $U=\operatorname{Lie}B$, and
$T=k\otimes_{\mathbf Z}\Lambda$, where $\Lambda$ is the
cocharacter lattice of the torus. Write its exact tangent sequence
as
\[
                 0\longrightarrow T\xrightarrow{j}V
                    \xrightarrow{P}U\longrightarrow0.
\]
Suppose $B$ has the self-dual line bundle of
Proposition~\ref{geo-proper-formal-hodge-comparison}, and that the
rigidified line bundle defining the formal module on $G$ is the
pullback of this line bundle. Give both biextensions their
rigidified adjunctions.
Then the module $V\oplus V^*$ with action
\eqref{geo-formal-hodge-action} is projective over $C$.
If its characteristic polynomials satisfy
the $r$-fold regular identity for an integer $r\geq0$,
it is free of rank $r$.
\end{proposition}
\begin{proof}
Let $e_G,e_B$ be the frames, and write their comparison as
$e_G(x,y)=\ell(x,y)^{-1}e_B(qx,qy)$, with $\ell$ equal to one
on both axes. The two induced adjunctions agree: their ratio is
a global unit on $G\times G$ equal to one on both axes, and hence
is one by the theorem on units on products
\cite[Theorem 1.1]{conrad-units-product}.
If $b_a,c_a$ are the mixed coefficients of the two adjunctions
and $Z$ is that of $\ell$, their relation is therefore
\[
 u_a^G(x,y)=u_a^B(qx,qy)\ell(x,\bar a y)/\ell(ax,y).
\]
Taking the mixed bilinear coefficient gives, for every $a$,
\begin{equation}\label{geo-toric-gauge-comparison}
 b_a+\rho(a)^tZ-Z\rho(\bar a)=P^t c_aP.
\end{equation}
The change $(v,\lambda)\mapsto(v,\lambda+Z^tv)$ identifies the
action with bilinear term $P^tc_aP$ with the original action.
In these new coordinates, the map
\[
       (v,\lambda)\longmapsto j^*\lambda
          \quad\hbox{from }V\oplus V^*
          \hbox{ to }T^\dagger
\]
is equivariant and surjective; the dagger denotes the
conjugate dual action. Its kernel is $V\oplus P^*U^*$.
The equivariant map
\[
           (v,P^*\mu)\longmapsto(Pv,\mu)
\]
from this kernel onto the proper abelian module $U\oplus U^*$
has kernel $T\oplus0$.
Thus the three successive quotients are
$T$, $U\oplus U^*$, and $T^\dagger$.

The $\mathcal O$-module $\Lambda$ is torsion-free: if $a\ne0$
kills an element, multiplying by an integral $b$ with
$ab=N_{E/\mathbf Q}(a)\ne0$ would give integer torsion.
It is therefore projective over the Dedekind domain $\mathcal O$.
The same argument applies to its integer dual, with conjugate
action. Their scalar extensions give the two projective toric
quotients. The middle quotient is projective by
Proposition~\ref{geo-proper-formal-hodge-comparison}.
Both exact sequences therefore split as $C$-modules, proving
projectivity of the full module. The residue-character argument
in Lemma~\ref{geo-witt-coefficient-freeness} gives the common rank.
No trace pairing is used at a ramified prime.
\end{proof}

These are statements about the prescribed formal module, rather
than its semisimplification. They pass to the actual local
lattice as follows. Let $S$ be a local coefficient ring and let
$\mathcal H$ be the underlying finite free $S$-module with its
$\mathcal O$-action. If its geometric closed fibre is free of rank
$r$ over $\mathcal O\otimes k$, then its residue-field fibre is
free of that rank: on each Artin local factor, choose $r$ generators
modulo the maximal ideal, apply Nakayama, and compare dimensions
after algebraic closure. Lift these generators to $\mathcal H$.
The resulting map
$(S\otimes_{\mathbf Z}\mathcal O)^r\to\mathcal H$ is a map between
free $S$-modules of the same rank and is an isomorphism modulo
the maximal ideal. Its determinant is a unit, so it is an
isomorphism over $S$.

For the boundary application, identify the prescribed integral
formal charts and their residue-fibre semiabelian quotients with
the original self-dual pullback line bundle. The propositions
then supply the frame comparison and coefficient freeness,
including at ramified primes. The next construction supplies
both filtrations on these same charts. The actual boundary
model and its complex multiplication evaluation remain separate.
